\section{Описание}

Требуется реализовать собственный словарь на базе красно-чёрного дерева. В основе лежат специфические правила раскраски узлов и операции балансировки — левое и правое вращение, а также перекраска, которые гарантируют, что высота дерева не превысит $2 \log (n + 1)$. 

Для этого реализуется класс \texttt{RBTree} и структура \texttt{Node}, включающая в структуру \texttt{Pair}, в которой хранится ключ, значение, указатели на потомков и структура данных \texttt{ParentWithColor}, которая хранит в себе ссылку на предка и цвет узла (RED/BLACK) с помощью техники \textit{Pointer tagging}, позволяя использовать младшие биты для хранения информации.

А после надо провести исследование скорости выполнения и потребления оперативной памяти.
\pagebreak

\subsection{Color}

Перечисление с жесткой типизацией (\texttt{enum class}), определяющее цвет узла. В качестве базового типа используется \texttt{bool}, что идеально подходит для последующей упаковки в младший бит указателя (0 для RED, 1 для BLACK).

\begin{table}[!ht]
\centering
\begin{tabular}{|p{7.5cm}|p{7.5cm}|}
\hline
\rowcolor{lightgray}
\multicolumn{2}{|c|}{Color} \\
\hline
\texttt{RED} & Константа, представляющая красный цвет узла \\
\hline
\texttt{BLACK} & Константа, представляющая черный цвет узла \\
\hline
\texttt{operator<<} & Перегрузка оператора вывода для отладки состояния дерева \\
\hline
\end{tabular}
\end{table}

\subsection{Pair}

Структура данных, представляющая собой полезную нагрузку узла. Она связывает строковый ключ (телефон или слово) с его соответствующим числовым значением \texttt{uint64\_t}.

\begin{table}[!ht]
\centering
\begin{tabular}{|p{7.5cm}|p{7.5cm}|}
\hline
\rowcolor{lightgray}
\multicolumn{2}{|c|}{Pair} \\
\hline
\texttt{std::string key} & Ключ словаря для лексикографического сравнения \\
\hline
\texttt{uint64\_t value} & Значение, ассоциированное с данным ключом \\
\hline
\texttt{Pair(...)} & Конструкторы для инициализации данных \\
\hline
\texttt{operator<<} & Форматированный вывод содержимого пары \\
\hline
\end{tabular}
\end{table}

\subsection{ParentWithColor}

Шаблонный класс, реализующий технику \textit{pointer tagging}. Поскольку адреса объектов в куче выровнены, младшие биты указателя всегда нулевые. Мы используем самый младший бит для хранения цвета, что экономит память и улучшает локальность данных.

\begin{table}[!ht]
\centering
\begin{tabular}{|p{7.5cm}|p{7.5cm}|}
\hline
\rowcolor{lightgray}
\multicolumn{2}{|c|}{ParentWithColor<T>} \\
\hline
\texttt{uintptr\_t data} & Числовое представление указателя, совмещенное с битом цвета \\
\hline
\texttt{struct ColorProxy} & Прокси-объект для удобного чтения и записи цвета через оператор \texttt{=} \\
\hline
\texttt{struct ParentProxy} & Прокси-объект для обращения к родителю и его разыменования \\
\hline
\texttt{color()} & Доступ к биту цвета узла \\
\hline
\texttt{parent()} & Доступ к указателю на родительский узел \\
\hline
\texttt{ptr() const} & Извлечение чистого адреса памяти без бита цвета \\
\hline
\end{tabular}
\end{table}

\subsection{Node}

Базовый узел дерева. Содержит данные, ссылки на детей и упакованный указатель на родителя. Использование \texttt{alignas} гарантирует корректное выравнивание для работы \textit{pointer tagging}.

\begin{table}[!ht]
\centering
\begin{tabular}{|p{7.5cm}|p{7.5cm}|}
\hline
\rowcolor{lightgray}
\multicolumn{2}{|c|}{Node} \\
\hline
\texttt{Pair data} & Полезная нагрузка узла \\
\hline
\texttt{Node* left / Node* right} & Указатели на левое и правое поддеревья \\
\hline
\texttt{ParentWithColor<Node> pwc} & Упакованный родитель и цвет узла \\
\hline
\texttt{Node(const Pair\&)} & Конструктор инициализации узла данными \\
\hline
\texttt{operator<<} & Вывод узла (для отладки: данные + цвет) \\
\hline
\end{tabular}
\end{table}

\subsection{RBTree}

Класс \texttt{RBTree} реализует самобалансирующееся бинарное дерево поиска. Все основные операции выполняются за $O(\log n)$ благодаря поддержанию свойств красно-чёрного дерева.

Для упрощения логики используется специальный фиктивный узел \texttt{NIL}, представляющий листья дерева. Он всегда окрашен в черный цвет и заменяет \texttt{nullptr}.

\begin{table}[!ht]
\centering
\begin{tabular}{|p{7.5cm}|p{7.5cm}|}
\hline
\rowcolor{lightgray}
\multicolumn{2}{|c|}{Поля класса} \\
\hline
\texttt{Node* root} & Корень дерева \\
\hline
\texttt{Node* NIL} & Фиктивный листовой узел (sentinel) \\
\hline
\end{tabular}
\end{table}

\begin{table}[!ht]
\centering
\begin{tabular}{|p{7.5cm}|p{7.5cm}|}
\hline
\rowcolor{lightgray}
\multicolumn{2}{|c|}{Методы модификации и поиска} \\
\hline
\texttt{bool insert(const Pair\& data)} & Вставка элемента с последующей балансировкой через \texttt{insertFixUp}. \\
\hline
\texttt{bool remove(const std::string\& key)} & Удаление элемента по ключу с возможной балансировкой через \texttt{removeFixUp}. \\
\hline
\texttt{Node* search(const std::string\& key)} & Поиск элемента (регистронезависимый). \\
\hline
\texttt{Node* search\_impl(const std::string\& key)} & Внутренняя реализация поиска (возвращает \texttt{NIL} при отсутствии). \\
\hline
\end{tabular}
\end{table}

\begin{table}[!ht]
\centering
\begin{tabular}{|p{7.5cm}|p{7.5cm}|}
\hline
\rowcolor{lightgray}
\multicolumn{2}{|c|}{Методы балансировки} \\
\hline
\texttt{void leftRotate(Node* x)} & Левый поворот вокруг узла \texttt{x}. \\
\hline
\texttt{void rightRotate(Node* y)} & Правый поворот вокруг узла \texttt{y}. \\
\hline
\texttt{void insertFixUp(Node* z)} & Восстановление свойств дерева после вставки. \\
\hline
\texttt{void removeFixUp(Node* x)} & Восстановление свойств дерева после удаления. \\
\hline
\end{tabular}
\end{table}

\begin{table}[!ht]
\centering
\begin{tabular}{|p{7.5cm}|p{7.5cm}|}
\hline
\rowcolor{lightgray}
\multicolumn{2}{|c|}{Служебные методы работы с деревом} \\
\hline
\texttt{Node* successor(Node* x)} & Поиск следующего по порядку узла. \\
\hline
\texttt{Node* minimum(Node* x)} & Поиск минимального элемента в поддереве. \\
\hline
\texttt{void clear(Node* node)} & Рекурсивное удаление поддерева. \\
\hline
\texttt{void setUpNIL()} & Инициализация фиктивного узла \texttt{NIL}. \\
\hline
\texttt{std::string to\_lower(std::string)} & Приведение строки к нижнему регистру. \\
\hline
\end{tabular}
\end{table}

\begin{table}[!ht]
\centering
\begin{tabular}{|p{7.5cm}|p{7.5cm}|}
\hline
\rowcolor{lightgray}
\multicolumn{2}{|c|}{Сериализация} \\
\hline
\texttt{void save(const std::string\& path)} & Сохранение дерева в бинарный файл. \\
\hline
\texttt{void load(const std::string\& path)} & Загрузка дерева из файла. \\
\hline
\texttt{void save\_impl(Node*, std::ofstream\&)} & Рекурсивная сериализация (pre-order). \\
\hline
\texttt{Node* load\_impl(std::istream\&, Node*)} & Рекурсивная десериализация дерева. \\
\hline
\end{tabular}
\end{table}

\pagebreak

\section{Консоль}

\begin{alltt}
root@74b1f15b9ece:/workspaces/Discrete-Analysis-MAI# cmake -B build
root@74b1f15b9ece:/workspaces/Discrete-Analysis-MAI# cmake --build build
root@74b1f15b9ece:/workspaces/Discrete-Analysis-MAI# ./build/lab2/lab2

+ a 1
OK
+ A 2
Exist
A
OK: 1
- A
OK
a
NoSuchWord

\end{alltt}

\pagebreak